% §VII --- Discussion. Status: REWRITTEN for the v8 admission-regime verdict.
\section{Discussion}\label{sec:discussion}

\subsection{Dual pricing is a domain property, not a method property}
The central methodological lesson of Section~\ref{sub:pathology} is that
the \emph{same} Lagrangian machinery, unchanged, exhibits both textbook
behaviours depending on where it operates. In the near-feasible nominal
domain the quality-critical multiplier climbs slowly and stays interior
($3.29$ at episode $500$, never touching its cap of $8$) against a realized
cost that closes on its limit---a shadow price forming on the constraint
boundary. In the peak overload domain the same multiplier ratchets
monotonically to whatever ceiling it is given ($\sim\!18.4$ under a cap of
$20$ across two earlier generations; exactly $8.00$ on $3/3$ seeds under
the v8 cap) against a cost stuck an order of magnitude above its limit.
This is not a tuning failure: CMDP theory says a bounded interior
multiplier exists only when the constraint set is
feasible~\cite{altman1999constrained,chow2017risk}. The pin \emph{is} the
optimizer correctly reporting ``no feasible policy exists here.''

Two experimental probes confirm the reading from opposite sides. The
$150$-episode dual warm-up lets the policy form before any constraint
price moves, yet once ascent resumes the cost is still $\approx\!0.37$ and
the climb recurs---shaping cannot manufacture a feasible region the
environment does not contain. Conversely, warm-\emph{starting} the
multiplier at its pinned value (Section~\ref{sub:m3}) changes nothing
either: the trajectory's destiny is set by the domain, not by its initial
condition or its schedule.

\subsection{In an infeasible domain, the cap is a behaviour knob---so use
admission, not price}
A subtler and, we believe, underreported effect: once the multiplier is
pinned, its \emph{ceiling} stops being a safety parameter and becomes a
behaviour-shaping knob. Capping $\lambda_{QC}$ at $8$ instead of $20$
halved the policy's critical-escalation rate ($0.126\to0.086$), dropping
it \emph{below} the certificate-derived budget band: at a pinned price of
$18.4$, escalation had been the cheap escape valve from the quality
penalty; at $8$ that pressure halves and the policy under-escalates.
Every choice of cap trades one criterion against another because the dual
no longer has an interior equilibrium to find. This coupling is exactly
why constraint \emph{pricing} is the wrong control surface for an
overloaded safety system, and why our architecture hands that regime to
the certified escalation/admission layer with an explicit budget
(Section~\ref{subsec:cert})---the dedicated feasibility mechanism that
reachability-constrained RL argues infeasible states
require~\cite{yu2022reachability}. The escalation-aware oracle makes the
point constructively: admitted-set mission success $0.916$ is achievable
at peak, but only by escalating $0.666$ of the critical load---no
strategy-layer price can be ``paid'' into a $0.155$ escalation budget when
the certificate says $0.105$ of the load is structurally unserveable and
the deadline geometry starves the rest.

\subsection{The constrained arm equalling the unconstrained arm is a
finding, not a defect}
At peak, the proposed constrained policy and its unconstrained ablation
converge to statistically indistinguishable behaviour. We report this
plainly because the mechanism is now visible: a saturated-and-uniform dual
penalty is a constant offset that cannot shape \emph{relative} action
values, so in the infeasible domain the dual term carries no gradient
information. The reward-wiring fix of Section~\ref{sub:fixes} is what
makes this interpretable---after the fix, the \emph{unconstrained} arm's
peak behaviour improved markedly (cache flooding $0.436\to0.298$, task
success $0.242\to0.384$), proving the reward path is load-bearing, while
the \emph{constrained} arm's peak behaviour stayed invariant across three
generations of reward and dual surgery (admitted mission
$0.179/0.179/0.180$). The invariance isolates the failure to the pricing
mechanism itself, not to plumbing. We spent three engineering generations
(v6 scale-consistency gates, v7 link/reward realism, v8 dual levers)
trying to make the peak dual price work before accepting the verdict; we
document that lineage as the motivation for the admission-regime framing
rather than presenting the architecture as designed-right-first-time.

\subsection{Communication as a flight-safety variable, corrected}
The safety bridge remains, but its sign discipline changed once the
coupling was corrected from substitution to queueing
(Sections~\ref{subsec:mg1} and~\ref{sub:bridge}): payload can only
\emph{add} waiting time to the tactical loop, so the honest headline is
asymmetric---heavy evidence \emph{taxes} the airspace ($+23.7$\,m,
$-6\%$ capacity at peak), light evidence \emph{avoids the tax}
($+9$\,cm), and nothing beats the BUBBLES baseline. The operational
co-design signal survives in this honest form: at low SNR the scheduler
should bias toward tokens not only for deadlines but to keep the
separation service at its certified baseline, precisely when spectrum is
scarcest. The claim is auditable at the parameter level against the
deliverable, including the contested C2-spectrum assumption, which we
bound by evaluating both shared and dedicated C2 modes.

\subsection{System closure with the companion study}
The companion paper~\cite{zhang2026askb4tx} establishes, in isolation, that
question-conditioned evidence routing beats full-image transmission---even
over an error-free channel---and that a zero-parameter rule matches a
data-calibrated selector. The present paper embeds that per-question
capability into a multi-UAV, safety-constrained scheduler and shows what
survives contention and regulation: the same evidence levels that win on
accuracy also keep the separation service at its baseline, and the
scheduling question becomes one of admission under overload. Together the
two studies form a closed loop from per-query evidence selection up to
airspace-level resource and safety management. A federated
extension---sharing evidence-selection statistics across UAVs and edge
sites while keeping the safety kernel local---is the natural next step
toward a deployable U-space semantic service.

\subsection{Limitations}\label{sub:limitations}
We keep the red lines explicit. \emph{First}, the peak condition is
reported as a stress test that exceeds the strategy layer's feasible
region: admitted-set mission success of the learned policy ($0.180$) sits
far below the escalation-aware oracle ($0.916$), and we claim only that
the safety accounting is carried structurally there---by the certificate
and the budgeted escalation channel---not that the learned policy closes
the gap. \emph{Second}, the certificate and the quality constraints
inherit the calibrated LUT quality model; a backend swap (legacy LUT)
shifts absolute numbers (the dual pin persists, $2/2$ seeds), so claims
are conditional on the calibration. \emph{Third}, the conflict-budget
normalization constant ($0.08$) is an engineering mapping within the
BUBBLES barrier-ladder framework, declared as a design choice rather than
a value read from the deliverable; the same holds for the M/G/1 offered-load
presets ($\rho$ values) behind Table~\ref{tab:wsens}. \emph{Fourth},
the escalation channel's budget is calibrated from the certificate roll,
but the $0.05$ headroom is a design margin; under the v8 cap the
learned policy under-escalates ($0.086$ vs.\ the $[0.105,0.205]$ band),
and the zero-shot blockage profile shows the learned shedding behaviour
does not transfer to unseen deeply-infeasible link regimes
(Section~\ref{sub:m4})---both reported as open behaviour-shaping issues
of pinned-dual regimes rather than hidden. \emph{Fifth}, per-seed variance is elevated by the
no-GAE Monte-Carlo advantage estimate, a deliberate simplicity trade for
the short dense-reward episodes. Finally, we do \emph{not} claim a
certified system: the framework is BUBBLES-conformant at the
parameter level, and all safety statements are simulation statements about
that model, not about flight-tested hardware.
